\section{Inherited physical interface and alias notation}
\label{sec:tpc39-interface}

We use the physical packet and exact CRT face calculus of TPC--37 and
TPC--38.  This section restates only the data needed to distinguish the
unconditional tomography from its conditional arithmetic application.

\subsection{Rows, targets, and the equality output}

Fix \(h\in\Z\setminus\{0\}\).  On a retained dyadic packet,
\begin{equation}
 Q=L_0D_0,
 \qquad QJ\asymp X,
 \label{eq:tpc39-interface-scales}
\end{equation}
where \(L_0\) and \(D_0\) denote the two source scales, so that \(L\) remains
reserved for the alias radius.  A physical row is
\begin{equation}
 \alpha=(\ell_\alpha,d_\alpha),
 \qquad
 \ell_\alpha\asymp L_0,
 \qquad
 d_\alpha\asymp D_0,
 \qquad
 m_\alpha=\ell_\alpha d_\alpha\asymp Q,
 \label{eq:tpc39-interface-row}
\end{equation}
and the source separation makes \(m_\alpha\) determine \(\alpha\).  The
two row coefficient systems obey
\begin{equation}
 \sum_\alpha|\gamma_\alpha^{(i)}|^2
 \ll X^{o(1)}Q,
 \qquad i=1,2.
 \label{eq:tpc39-interface-row-l2}
\end{equation}
For \(j\) in a positive interval \(\cJ\) of \(O(J)\) integers, put
\begin{equation}
 N_\alpha(j)=m_\alpha j+h\asymp X.
 \label{eq:tpc39-interface-target}
\end{equation}

The retained output space is
\begin{equation}
 \cH_{\rm phys}
 =\ell^2\{(L,\gamma,j)\}
 \oplus\ell^2\{(R,\alpha,j)\}.
 \label{eq:tpc39-interface-Hphys}
\end{equation}
For a determinant shift \(F\), let
\(\mathscr Z_F[\kappa]\in\cH_{\rm phys}\) denote the shifted output on
\begin{equation}
 su=N_\alpha(j)+F,
 \label{eq:tpc39-interface-shifted-fiber}
\end{equation}
with target atom multiplier \(\kappa_N(F)\).  Write
\begin{equation}
 \mathscr R_F:=\mathscr Z_F[\one]
 \label{eq:tpc39-interface-raw-column}
\end{equation}
for the raw column.  The original regular--subpacket output is
\begin{equation}
 \mathscr Y_{\rm phys}=\mathscr R_0,
 \qquad
 \|\mathscr R_0\|^2=V_L+V_R.
 \label{eq:tpc39-interface-physical-gate}
\end{equation}
The inherited target is
\begin{equation}
 \boxed{
 \|\mathscr R_0\|^2
 \ll X^{o(1)}\frac{Q^3}{J}.}
 \label{eq:tpc39-interface-target-gate}
\end{equation}

\subsection{The three--modulus completion}

Choose distinct primes \(q_1,q_2,q_3\asymp J\) satisfying the inherited
range conditions
\begin{equation}
 \max\cJ<q_i,
 \qquad
 \max_\alpha d_\alpha<q_i<\min_\alpha\ell_\alpha,
 \qquad q_i\nmid h,
 \label{eq:tpc39-interface-prime-range}
\end{equation}
and put
\begin{equation}
 M=q_1q_2q_3.
 \label{eq:tpc39-interface-M}
\end{equation}
TPC--37 absorbs the row--time pairs for which
\((N_\alpha(j),M)>1\) at the inherited acceptable scale.  All physical
outputs below refer to the complementary regular subpacket.  No claim is
made that the additive twists introduced later supply this zero--face
estimate; it is inherited before completion.

TPC--38 constructs the full and proper target fields \(C_N(F)\) and
\(P_N(F)\).  Their exact recombination is
\begin{equation}
 \boxed{C_N(F)+P_N(F)=\one_{M\mid F}.}
 \label{eq:tpc39-interface-recombination}
\end{equation}
Let \(B\) be an integer such that every determinant used by the completed
packet satisfies \(|F|\le B\).  We assume \(B\ge M\), equivalently
\(L\ge1\), as is automatic in the endpoint regime considered below, and put
\begin{equation}
 L=\left\lfloor\frac BM\right\rfloor,
 \qquad
 \cI_B=\{-L,\ldots,L\},
 \qquad
 K_B=2L+1.
 \label{eq:tpc39-interface-alias-ledger}
\end{equation}
For \(k\in\cI_B\), define
\begin{equation}
 \mathscr Z_k:=\mathscr R_{kM}.
 \label{eq:tpc39-interface-Zk}
\end{equation}
Summing \eqref{eq:tpc39-interface-recombination} before estimating gives
\begin{equation}
 \sum_{|F|\le B}
 \bigl(\mathscr Z_F[C]+\mathscr Z_F[P]\bigr)
 =\sum_{k=-L}^{L}\mathscr Z_k.
 \label{eq:tpc39-interface-completed-alias}
\end{equation}
The right side is a congruence completion, not \(\mathscr Z_0\).

\subsection{Endpoint normalization}

At the endpoint,
\begin{equation}
 Q=X^{267/400+o(1)},
 \qquad
 J=X^{133/400+o(1)},
 \qquad
 M=X^{399/400+o(1)}.
 \label{eq:tpc39-interface-endpoint}
\end{equation}
For \(B\asymp X\),
\begin{equation}
 K_B\asymp 1+\frac XM
 =X^{1/400+o(1)}.
 \label{eq:tpc39-interface-K}
\end{equation}
Put
\begin{equation}
 B_0=Q^2J.
 \label{eq:tpc39-interface-B0}
\end{equation}
Then
\begin{equation}
 \boxed{
 B_0K_B
 =X^{o(1)}\frac{Q^3}{J}.}
 \label{eq:tpc39-interface-key-ledger}
\end{equation}
The proved input--copy diagonal has size \(X^{o(1)}B_0\), but it is not a
bound for the raw output column \(\mathscr Z_0\).  The latter remains the
physical gate in \eqref{eq:tpc39-interface-target-gate}.
